% Page budget: 1.55 pages | Owns: augmentation topology, additional states, encoder, and closed-loop objective.
% WRITING GUIDE
% Purpose: Define the final augmentation architecture and explain how it is trained inside the known feedback loop.
% Start here: Current implementation decisions D-129, D-140, D-141, D-180, D-220, D-233, and D-237 in docs/decisions.md.
% Supporting material: Quinten's 18 September outline feedback, the Meeting-audit synthesis, Thesis-writeup/Documentation/TRAINING-DESIGN.md, and docs/writeup figures.
% Mandatory papers: Read Hoekstra's 2025 LFR augmentation paper, the 2026 first-principles augmentation paper, the encoder-initialization paper, and Kessels' Chapter 5 before writing this section.
% Research-plan anchor: Preserve Aspect 2's dynamic parallel augmentation objective; document the departure from open-loop training and earlier state routing.
% Current decisions: Separate physical and additional states, keep the controller outside the model, score the complete training window without burn-in, and use one network for the physical correction and added-state update.
% Choices to justify: Final reported routing, added-state dimension, ANN width and depth, zero initialization, encoder history, closed-loop objective, rollout length, full-window scoring, and validation horizon. Use the configuration of the reported thesis run, not a superseded routing experiment.
% Horizon check: Treat encoder history, scored training window, joint-estimation horizon, and free-run validation horizon separately; relate candidates to relevant stable modes and test sensitivity to slower dynamics and free-integrator accumulation.
% Implementation audit: Read scripts/gantry/gantry_interconnect_dynamic.py, gantry_dynamic/{model,config,controller,data,training}.py, and model_augmentation/fit_systems/{interconnect,closed_loop,pre_encoder}.py.
% Reference check: Verify sources for parallel LFR augmentation, encoder initialization, closed-loop state extension, and any claimed architectural precedent against the original paper or thesis.
% Cautions: Earlier routing plans are superseded; distinguish the truth-only hidden absorber from learned states; closed-loop fit alone does not prove autonomous plant recovery.
% Figures: Absorber schematic, author's updated architecture, and thesis-owned common-reference loops. No residual panel or horizon figure.
% Section is finished when: Every signal has a defined frame and timing, the RK4 boundary is explicit, and the described architecture matches the final code.
%
% SOURCE MAP
% Hoekstra et al. 2025 EJC: dynamic-parallel topology, additional states, encoder-based truncated simulation.
% Hoekstra et al. 2026 Automatica submission: extended LFR formulation, baseline-preserving model/encoder initialization, joint identification context.
% Hoekstra et al. 2026 encoder paper: reconstructability-based W^b initialization; it does not identify W^a or test dynamic augmentation.
% Kessels 2025 thesis: direct closed-loop rollout, controller equations, and reconstruction of the machine controller state for each window.
% Drenth 2025 thesis: LPV-specific augmentation precedent; not the unrestricted MLP structure used here.
% This project: gantry specialization, residual-loop implementation, exact routing, RK4 boundary, and verification.
\section{Dynamic LPV-LFR Augmentation}\label{sec:augmentation}

The baseline of Section~\ref{sec:lfr} describes the rigid-body dynamics of
the gantry, but not every effect of the machine.
Garc\'ia-Herreros et al.\ neglect the resonance of the supporting base, the
vibration of the cross-arm on its flexible plates, the detent forces of the
linear motors, and the variation of friction along the
stroke~\cite[Sec.~2.4]{garcia2013model}, and our baseline \eqref{eq:eom}
also leaves out the Coulomb friction that their model contains.
This section augments the baseline with a learned component for such
effects, and estimates the physical parameters $\theta_{\mathrm{base}}$, the
network parameters $\theta_{\mathrm{aug}}$ and the encoder parameters $\eta$
jointly from records measured under feedback.
Section~\ref{sec:aug_structure} defines the augmented model,
Section~\ref{sec:aug_encoder} the encoder that supplies the initial state of
each training window, Section~\ref{sec:aug_closed_loop} the closed-loop
rollout and the objective, and Section~\ref{sec:aug_init} the initialisation
from the baseline.

\subsection{Augmented model}\label{sec:aug_structure}
Several of the omitted effects are dynamic: a vibration mode of the base or
the cross-arm adds poles that the six states of the baseline cannot
represent.
Refitting $\theta_{\mathrm{base}}$ cannot add this model order, and neither
can a static correction of the state update, because both only change the
dynamics of the existing states~\cite[Sec.~2]{hoekstra2025lfr}.
We therefore add states to the model and identify the dynamic-parallel
structure~\cite[Table~1]{hoekstra2026lfr}
\begin{equation}
\begin{aligned}
 \tilde x_{k+1}
 &=f_{\mathrm{base}}(\theta_{\mathrm{base}},\tilde x_k,u_k)
   +f_{\mathrm{aug}}(\theta_{\mathrm{aug}},\hat x_k,u_k),\\
 \bar x_{k+1}
 &=g_{\mathrm{aug}}(\theta_{\mathrm{aug}},\hat x_k,u_k),\\
 \hat y_k&=h_{\mathrm{base}}(\tilde x_k),
\end{aligned}
\label{eq:aug_transition}
\end{equation}
where $\tilde x=\col(q,\dot q)\in\R^6$ is the physical state of
Section~\ref{sec:lfr}, $\bar x\in\R^{n_{x_a}}$ collects the additional
states, $\hat x=\col(\tilde x,\bar x)$, $u$ is the generalised force of
\eqref{eq:Ptransform}, and $\hat y$ predicts the measured positions
$y=\qact$.
The baseline transition $f_{\mathrm{base}}$ keeps its physical parameters,
and the learned functions $f_{\mathrm{aug}}$ and $g_{\mathrm{aug}}$ act in
parallel to it, so $\theta_{\mathrm{base}}$ keeps the meaning it has in
\eqref{eq:eom}.
Drenth formulates this augmentation, additional states included, for an
LPV-LFR baseline~\cite[Sec.~5.2]{drenth2025thesis}.

The baseline transition $f_{\mathrm{base}}$ is one step of the classical
fourth-order Runge-Kutta (RK4) method, applied to the state equation of
\eqref{eq:eom}, written as $\dot{\tilde x}=F(\theta_{\mathrm{base}},\tilde
x,u)$, over one sample period $T_s$.
The input is held over the step, as the actuator force is held between
samples.
Each of the four stages evaluates $M(Y)$ at its own intermediate state, so
the scheduling variable moves within the step and $f_{\mathrm{base}}$
remains self-scheduled.
The learned correction is added after the step: $f_{\mathrm{base}}$
integrates the continuous-time physics, whereas $f_{\mathrm{aug}}$ is a
discrete-time correction of the sampled state.

The two learned functions are the outputs of one multilayer perceptron
$\phi_{\mathrm{aug}}$ with tanh hidden layers and a linear output layer,
\begin{equation}
 \begin{bmatrix}
  f_{\mathrm{aug}}(\theta_{\mathrm{aug}},\hat x_k,u_k)\\
  g_{\mathrm{aug}}(\theta_{\mathrm{aug}},\hat x_k,u_k)
 \end{bmatrix}
 =\phi_{\mathrm{aug}}(\theta_{\mathrm{aug}},\hat x_k,u_k),
 \label{eq:aug_mlp}
\end{equation}
where $\theta_{\mathrm{aug}}$ collects its weights and biases.
The first six outputs form $f_{\mathrm{aug}}$ and the last $n_{x_a}$ form
$g_{\mathrm{aug}}$, so the update of the additional states and the
correction of the physical states share one network
(Fig.~\ref{fig:aug_structure}).
The correction $f_{\mathrm{aug}}$ writes to all six physical rows, the
positions included.
The omitted effects listed above include forces on every motor, such as
detent and friction forces, so the correction must reach all three axes.
Within one sample, a force held over the step changes the position as well
as the velocity, because the input-to-state map of a held input has nonzero
position rows; a discrete-time correction of such a force therefore needs
the position rows.

\begin{figure}[t]
 \centering
 \resizebox{\columnwidth}{!}{\input{figures/tex/augmentation_structure}}
 \caption{Augmented model. The physics step $f_{\mathrm{base}}$ and the
 network $\phi_{\mathrm{aug}}$ act in parallel on $\hat x_k$;
 $\phi_{\mathrm{aug}}$ supplies $f_{\mathrm{aug}}$ and $g_{\mathrm{aug}}$,
 and the output map is not learned ($h_{\mathrm{aug}}=0$). The encoder
 $\psi$ sets $\hat x$ at each window start. During training, $u_k$ is the
 closed-loop input $\hat u_k$ of Section~\ref{sec:aug_closed_loop}.}
 \label{fig:aug_structure}
\end{figure}

The output map $h_{\mathrm{base}}(\tilde x)=\begin{bmatrix}P^\top &
0\end{bmatrix}\tilde x$ is the kinematic map of \eqref{eq:Ptransform}, and it
is not augmented ($h_{\mathrm{aug}}=0$), so the measured positions are read
from the physical state alone.
Within one step, no learned function feeds the input of another, so the
computational graph of the model is acyclic.
The state equation $F$ is smooth because $M(Y)$ is invertible for every
iterate of $\theta_{\mathrm{base}}$ (Section~\ref{sec:params}), and together
with the tanh activations and the acyclic graph this makes the augmented
model well posed~\cite[Thm.~9]{hoekstra2026lfr}.
The additional states have no assigned physical meaning and reach the output
only through $f_{\mathrm{aug}}$.
In particular, they are not the states of the absorber that generates the
simulated data of Section~\ref{sec:setup}.

\subsection{Encoder}\label{sec:aug_encoder}
Each training window simulates \eqref{eq:aug_transition} from an initial
state at the window start $\tau$, but only the positions are measured.
Following Beintema et al.~\cite{beintema2023subnet}, an encoder estimates
this state from the measured input and output history.
We use the encoder structure of~\cite[Eq.~(8)]{hoekstra2026encoder}, a
linear map with a nonlinear correction,
\begin{equation}
\begin{aligned}
 \hat x_{\tau|\tau}
 &=\psi_\eta(\xi_\tau)
 =\begin{bmatrix}W^b\\ W^a\end{bmatrix}\xi_\tau+\tilde\psi(\xi_\tau),\\
 \xi_\tau&=\col\bigl(y_{\tau-n_a:\tau},\,u_{\tau-n_b:\tau}\bigr),
\end{aligned}
 \label{eq:aug_encoder}
\end{equation}
where $W^b$ gives the physical states, $W^a$ the additional states,
$\tilde\psi$ is an MLP, and $\eta$ collects $W^b$, $W^a$ and the parameters
of $\tilde\psi$.
Both histories include the current sample $\tau$, as in the
reconstructability map of~\cite[Eq.~(15)]{hoekstra2026encoder} from which
Section~\ref{sec:aug_init} initialises $W^b$.
The encoder estimates every state, the measured positions included, and is
trained jointly with the model through the objective of
Section~\ref{sec:aug_closed_loop}.

\subsection{Closed-loop rollout and objective}\label{sec:aug_closed_loop}
The records are measured under feedback, and the model is meant to predict
the machine under feedback: the response to a given reference, feedforward
and known controller, also when that controller is changed.
An open-loop simulation driven by the recorded input, as in the truncated
objective of~\cite[Eq.~(22)]{hoekstra2026lfr}, does not reproduce this use.
It also lets errors accumulate on the gantry: $K$ in
\eqref{eq:damping_stiffness_matrices} has no stiffness on $X$ and $Y$, so a
velocity error integrates into a position error that persists in an
open-loop rollout.
Inside the loop, the controller acts on such an error as it does on the
machine, provided the model in feedback with the controller is stable.
Kessels trains an augmented model with the known feedback controller inside
each truncated rollout and propagates gradients through
it~\cite[Eqs.~(5.12), (5.13c), (5.13d)]{kessels2025ai}.
We adopt this rollout and keep the controller outside the learned model, so
that it can be replaced in evaluation, as Kessels does for a changed
controller~\cite[Sec.~5.3.2.3]{kessels2025ai}.

The recorded loop is driven by a reference $r$ and a feedforward
$u^{\mathrm{ff}}$ through a linear controller with state $s$ and matrices
$(A_K,B_K,C_K,D_K)$, which maps the servo error to actuator forces.
The model loop receives the same $r$, $u^{\mathrm{ff}}$ and controller, and
differs only in its output $\hat y$ (Fig.~\ref{fig:aug_closed_loop}):
\begin{equation}
\begin{aligned}
 \hat s_{k+1}&=A_K \hat s_k+B_K(r_k-\hat y_k),\\
 \hat u_{\mathrm{act},k}&=C_K \hat s_k+D_K(r_k-\hat y_k)+u^{\mathrm{ff}}_k,
 \qquad \hat u_k=P\hat u_{\mathrm{act},k}.
\end{aligned}
\label{eq:aug_direct_controller}
\end{equation}
The recorded loop satisfies the same equations with $s$, $y$ and $\uact$ in
place of $\hat s$, $\hat y$ and $\hat u_{\mathrm{act}}$.
Subtracting the recorded loop from the model loop removes $r$ and
$u^{\mathrm{ff}}$,
\begin{equation}
\begin{aligned}
 x^c_{k+1}&=A_K x^c_k+B_K(y_k-\hat y_k),\qquad x^c_\tau=0,\\
 \hat u_{\mathrm{act},k}&=u_{\mathrm{act},k}+C_K x^c_k+D_K(y_k-\hat y_k),
\end{aligned}
\label{eq:aug_controller}
\end{equation}
where $x^c=\hat s-s$ is the difference of the controller states.
The model thus receives the recorded input plus the controller's response to
its own output error.
This form is exact when both loops use the same linear controller, not
scheduled on the model state, with the same feedforward and without
saturation.

\begin{figure}[t]
 \centering
 \includegraphics[width=\columnwidth]{closed_loop_same_reference.pdf}
 \caption{Recorded plant loop (top) and augmented-model loop (bottom),
 driven by one shared reference $r_k$ and feedforward $u_k^{\mathrm{ff}}$,
 with the same feedback controller. The model output is evaluated
 before the input advances its state to the next sample.}
 \label{fig:aug_closed_loop}
\end{figure}

The condition $x^c_\tau=0$ starts the model's controller in the state of the
recorded controller, $\hat s_\tau=s_\tau$, without reconstructing that
state.
Kessels obtains the same initial condition by reconstructing the controller
state from the measured output, the reference and the
controller~\cite[Remark~5.4]{kessels2025ai}.
Model error before $\tau$ is therefore not carried into the window.
Because the model has no feedthrough from $u$ to $\hat y$, each sample is
evaluated in a fixed order: $\hat y_k$ from the state, then $\hat u_k$ from
\eqref{eq:aug_controller}, then both states advance.
The controller attenuates the effect of model error inside its bandwidth, so
a good closed-loop fit does not by itself show that the open-loop plant
model is accurate.

Over the set of window starts $\mathcal T$ and the window length $n_f$, the
parameters $\vartheta=(\theta_{\mathrm{base}},\theta_{\mathrm{aug}},\eta)$
minimise the output error
\begin{equation}
 V(\vartheta)=\frac{1}{|\mathcal T|\,n_f}
 \sum_{\tau\in\mathcal T}\sum_{\ell=0}^{n_f-1}
 \norm{T_y\bigl(y_{\tau+\ell}-\hat y_{\tau+\ell|\tau}\bigr)}_2^2
 \label{eq:aug_loss}
\end{equation}
subject to \eqref{eq:aug_transition}, \eqref{eq:aug_encoder} and
\eqref{eq:aug_controller}, where $T_y$ is the diagonal matrix of inverse
output standard deviations over the training records.
This is the truncated objective of~\cite[Eq.~(22)]{hoekstra2026lfr}, with
each window simulated in closed loop as in~\cite[Eq.~(5.12)]{kessels2025ai}.
As in~\cite[Sec.~5.3]{hoekstra2026lfr}, inputs, outputs and states are
scaled to zero mean and unit standard deviation over the training records,
with $f_{\mathrm{base}}$ evaluated inside the same scaling; this changes the
coordinates but not the dynamics, so the equations of this section are
written in physical units and the scaling enters the objective only through
$T_y$, which takes the place of Kessels' output weighting matrix.

The physical parameters, the network and the encoder are estimated jointly,
while the controller stays fixed.
The parameters $\theta_{\mathrm{base}}$ are the ten identifiable
combinations of Section~\ref{sec:params}, estimated in the bounded
coordinates defined there.
No parameter regularisation such as~\cite[Eq.~(26)]{hoekstra2026lfr} is
added: it limits the drift of $\theta_{\mathrm{base}}$, but it does not
remove the directions in which the network can cancel a parameter change,
which Section~\ref{sec:obc} addresses by construction.
Every sample of each window is scored.
The encoder history $n_a=n_b$, the window length $n_f$, the spacing of the
window starts and the free-run validation horizon are separate quantities;
Section~\ref{sec:setup} gives their values.

\subsection{Initialisation from the baseline}\label{sec:aug_init}
A random initialisation of the network can make the augmented model
unstable and the optimisation poor~\cite[Sec.~5.4]{hoekstra2026lfr}, whereas
a negligible initial network contribution makes an unstable closed-loop
model less likely as long as the controller stabilises the
baseline~\cite[Sec.~5.2.3]{kessels2025ai}.
Training therefore starts from the baseline.
The weights and biases of the output layer of $\phi_{\mathrm{aug}}$ are set
to zero, so $f_{\mathrm{aug}}=0$ and $g_{\mathrm{aug}}=0$ for every input,
and the augmented model reproduces the baseline from the same physical
initial state, as required by~\cite[Eq.~(29)]{hoekstra2026lfr}.
From this start, $f_{\mathrm{aug}}$ can also learn changes that a different
$\theta_{\mathrm{base}}$ would produce, so the split between the two
components is not unique~\cite[Sec.~5.2]{hoekstra2026lfr}.
Section~\ref{sec:obc} constrains this split.

The physical rows $W^b$ of the encoder are initialised so that the encoder
reconstructs the state of the baseline.
The reconstructability map is available in closed form for a linear model,
so we linearise \eqref{eq:eom} at the nominal parameters about rest at $Y=0$,
as proposed for nonlinear baselines in~\cite[Sec.~3.1]{hoekstra2026encoder}.
This gives
\begin{equation}
 A_c(Y)=\begin{bmatrix}0 & I\\ -M(Y)^{-1}K & -M(Y)^{-1}C\end{bmatrix},\qquad
 B_c(Y)=\begin{bmatrix}0\\ M(Y)^{-1}\end{bmatrix},
\label{eq:aug_lin}
\end{equation}
which we discretise at $Y=0$ by a zero-order hold over $T_s$, matching the
held input of the RK4 step, to obtain $(A_d,B_d)$ with output matrix
$C_d=\begin{bmatrix}P^\top & 0\end{bmatrix}$.
For this model, $W^b$ is the reconstructability
map~\cite[Eqs.~(16), (17)]{hoekstra2026encoder}, which exists because the
measured positions make $(A_d,C_d)$ observable, so the observability matrix
over the encoder history has full column rank.
The linearisation only sets the starting value of $W^b$, since the encoder
is trained jointly with the model.
The rows $W^a$ have no baseline counterpart and are drawn randomly, and the
output layer of $\tilde\psi$ starts at zero, so the encoder initially
returns its linear part.
Because $f_{\mathrm{aug}}=0$ at initialisation, the encoded additional
states do not affect the output of the initial model.
